\documentclass[11pt,a4paper,twoside,bibliography=totoc,listof=totoc,titlepage=true,chapterprefix=true]{scrreprt}

%% ---------------------------------------------------------------------------
%% Packages
%% ---------------------------------------------------------------------------

\usepackage[english]{babel} %Option: german für deutsche Variante
\usepackage[T1]{fontenc}   
\usepackage[utf8]{inputenc} 
\usepackage{lmodern}

\KOMAoptions{parskip=half, numbers=noenddot} 

\usepackage{amsmath}
\usepackage{amssymb}
\usepackage{graphicx}
\usepackage{subfig} % uses for subfloat: two pictues side by side
\usepackage{amsthm} % needed for proofs
\usepackage{multirow} % multirows for tables
\usepackage{placeins} % for floatbarrier
\usepackage{booktabs} % Für schönere wissenschaftliche Tabellen
\usepackage{tikz}
\usepackage[backend=biber, style=alphabetic]{biblatex} 
\usepackage[pdfpagelabels,hidelinks]{hyperref}
\usepackage{csquotes} % Wenn babel mit biblatex, dann wird csqoutes empfohlen, damit zitierter Text richtiges Typenset hat
\usepackage[numbered]{bookmark} % Für nummerierte PDF-Lesezeichen entsprechend der Chapter/Section-Nummer

\usepackage{geometry} % Einstellungen der Ränder 
\geometry{paperwidth=210mm, paperheight=297mm, height=247mm,width=167mm,top=21.5mm,right=15mm}%,headsep=7mm,headheight=10mm,footskip=11mm}
%\geometry{includehead,includefoot,inner=3cm,outer=2cm,top=2cm,bottom=2.5cm}

\usepackage{supertabular}
\usepackage{color} % Für farbige Boxen --> Gut für Notizen

\usepackage{xifthen}							% provides \isempty test

\usepackage[acronym,nonumberlist,section]{glossaries} % test this package

%\usepackage[ruled, vlined, linesnumbered,german]{algorithm2e}
\usepackage[ruled, vlined, linesnumbered]{algorithm2e}
\usepackage{listings}
\usepackage[headsepline, plainfootsepline]{scrlayer-scrpage}

%% ---------------------------------------------------------------------------
%% Definitions
%% ---------------------------------------------------------------------------

\newtheorem{theorem}{Theorem}[chapter]
\newtheorem{assumption}{Assumption}[chapter]

%% Heading styles
\RedeclareSectionCommand[%
		tocentryformat={\rmfamily\bfseries},%
		tocpagenumberformat=\rmfamily
	]{chapter}

\addtokomafont{part}{\rmfamily}
\addtokomafont{chapterprefix}{\huge\rmfamily}
\addtokomafont{chapter}{\Huge\rmfamily}
\addtokomafont{section}{\rmfamily}
\addtokomafont{section}{\rmfamily}
\addtokomafont{subsection}{\rmfamily}
\addtokomafont{subsubsection}{\rmfamily}
\addtokomafont{paragraph}{\rmfamily}
\addtokomafont{subparagraph}{\rmfamily}

%% Bibliography Styles
\renewcommand{\mkbibnamefirst}[1]{\textsc{#1}}
\renewcommand{\mkbibnamelast}[1]{\textsc{#1}}
\renewcommand{\mkbibnameprefix}[1]{\textsc{#1}}
\renewcommand{\mkbibnameaffix}[1]{\textsc{#1}}

%% Style for Algorithms
\newcommand{\listofalgorithmes}{\tocfile{\listalgorithmcfname}{loa}}
\SetKwRepeat{Do}{do}{while}%
\SetAlFnt{\small} %\SetAlFnt{\footnotesize }
\SetKwComment{Comment}{$\triangleright$\ }{}
\SetVlineSkip{2pt} 
\DontPrintSemicolon
\SetNlSty{}{}{:} % No Bold Line Numbers
\SetAlgoNlRelativeSize{-1} % relative Size of Line Number
\SetCommentSty{} % Sonst wäre Kommentar Fett    

%colorbox with line break
%background: colorbox adds some space before and after the text inside it, this space is \fboxsep
\newcommand{\hly}[1]{\colorbox{yellow}{\parbox{\dimexpr\textwidth-2\fboxsep\relax}{#1}}} 
\newcommand{\hlg}[1]{\colorbox{green}{\parbox{\dimexpr\textwidth-2\fboxsep\relax}{#1}}}
\newcommand{\hlr}[1]{\colorbox{red}{\parbox{\dimexpr\textwidth-2\fboxsep\relax}{#1}}}

%$ Listing Style
%% Listing Style for Matlab
\definecolor{mygreen}{rgb}{0,0.6,0}
\definecolor{mygray}{rgb}{0.5,0.5,0.5}
\definecolor{mymauve}{rgb}{0.58,0,0.82}
\lstset{numbers=left, numberstyle=\tiny, numbersep=5pt} \lstset{language=Matlab}
\lstset{basicstyle=\ttfamily}
\lstset{literate=%
	{Ö}{{\"O}}1
	{Ä}{{\"A}}1
	{Ü}{{\"U}}1
	{ß}{{\ss}}2
	{ü}{{\"u}}1
	{ä}{{\"a}}1
	{ö}{{\"o}}1
} 
\lstset{ %
	backgroundcolor=\color{white},   		% choose the background color; you must add \usepackage{color} or \usepackage{xcolor}
	basicstyle=\footnotesize,        					% the size of the fonts that are used for the code
	breakatwhitespace=false,         			% sets if automatic breaks should only happen at whitespace
	breaklines=true,                 								% sets automatic line breaking
	captionpos=b,                    								% sets the caption-position to bottom
	commentstyle=\color{mygreen},    	% comment style
	deletekeywords={...},            						% if you want to delete keywords from the given language
	escapeinside={\%*}{*)},          					% if you want to add LaTeX within your code
	extendedchars=true,              						% lets you use non-ASCII characters; for 8-bits encodings only, does not work with UTF-8
	frame=single,                    								% adds a frame around the code
	keepspaces=true,                 						% keeps spaces in text, useful for keeping indentation of code (possibly needs columns=flexible)
	keywordstyle=\color{blue},       				% keyword style
	language=Octave,                 						% the language of the code
	morekeywords={*,...},            						% if you want to add more keywords to the set
	numbers=left,                    									% where to put the line-numbers; possible values are (none, left, right)
	numbersep=5pt,                   						% how far the line-numbers are from the code
	numberstyle=\tiny\color{mygray}, 		% the style that is used for the line-numbers
	rulecolor=\color{black},         						% if not set, the frame-color may be changed on line-breaks within not-black text (e.g. comments (green here))
	showspaces=false,                						% show spaces everywhere adding particular underscores; it overrides 'showstringspaces'
	showstringspaces=false,          				% underline spaces within strings only
	showtabs=false,                  							% show tabs within strings adding particular underscores
	stepnumber=2,                    							% the step between two line-numbers. If it's 1, each line will be numbered
	stringstyle=\color{mymauve},     			% string literal style
	tabsize=2,                       										% sets default tabsize to 2 spaces
	title=\lstname                   									% show the filename of files included with \lstinputlisting; also try caption instead of title
}

% Kopf- und Fußzeile Einstellungen
\clearpairofpagestyles % Kopf und Fußzeile löschen
\pagestyle{scrheadings}
\lehead[]{\pagemark}
\rehead[]{\upshape{\headmark}}
\lohead[]{\upshape{\headmark}}
\rohead[]{\pagemark}
\automark[section]{chapter} % Kopfzeile: Linke Seite: Section, Rechte Seite: Chapter

\makeglossaries

\newcommand{\addWithPreCommaSpace}[1]{%
	\if\relax #1\relax% If statement for checking empty parameter
	% if #1 is empty
	\else%
	, #1% if #1 is not empty
	\fi%
}

%% ---------------------------------------------------------------------------
%% Thesis Information
%% ---------------------------------------------------------------------------

\newcommand{\Dachorganisation}{Karlsruher Institut für Technologie - KIT}
\newcommand{\Institut}{Institut für Regelungs- und Steuerungssysteme - IRS}
\newcommand{\titelderarbeit}{Title}
\newcommand{\nummerderarbeit}{Bxxx} % Or: Mxxx
\newcommand{\artderarbeit}{Bachelor's Thesis} % Or: Master's Thesis
\newcommand{\student}{Name Surname}
\newcommand{\studentgrad}{} % Or: B.Sc.
\newcommand{\betreuer}{Name Surname}
\newcommand{\betreuergrad}{M.Sc.}
\newcommand{\ausgabe}{10th April 2019}
\newcommand{\abgabe}{10th January 2020}
\newcommand{\sperrvermerkUnternehmen}{Xxxx Gmbh}
\newcommand{\sperrvermerkDatum}{10th January 2025} % 5 years after submission

%% ---------------------------------------------------------------------------
%% Additional Files and Setup
%% ---------------------------------------------------------------------------

\hypersetup{
	pdfpagelayout=TwoPageRight,
	pdfstartview=Fit,
	pdfauthor={\student},
	pdftitle={\titelderarbeit},
	pdfsubject={\artderarbeit}
}

\addbibresource{04_bib/bib.bib}
% Setup for Title Page
% titlepage

\renewcommand\maketitle{
\begin{titlepage}
%\sffamily
\newgeometry{left=20mm, right=15mm, top=25mm, bottom=30mm}
\hspace{0.5cm}
\includegraphics[width=0.95\textwidth]{03_images/title_page/kitlogo_fakultaet.pdf}


\rule{\textwidth}{0.6mm}

\vspace{2cm}
\hspace{0.5cm}
{\huge \textsf{\student}}



\vspace{6.25cm}

\hspace{0.5cm}\begin{minipage}{0.95\linewidth}
\raggedright

{\fontsize{30pt}{36pt}\selectfont \textsf{\textbf{\titelderarbeit}} \par}


\end{minipage}

\vspace{6.25cm}
\rule{\textwidth}{0.6mm}

\vspace{0.75cm}

\hspace{0.5cm} 
{\Large \textsf{ \textbf{\artderarbeit{}\ \nummerderarbeit}}}

\vspace{0.75cm}

\hspace{0.5cm}
\includegraphics[width=0.9\textwidth]{03_images/title_page/irslogo.pdf}

\vspace{0.5cm}


\rule{\textwidth}{0.6mm}



\end{titlepage}
}

% Start Dokument 
\begin{document}
	\pagenumbering{roman}
	
	%% ---------------------------------------------------------------------------
	%% ---------------------------------------------------------------------------
	
	\pdfbookmark{Front Page}{title}
	\maketitle % Command redefined in title_page.tex
	
	% If no non-disclosure notice is required, remove the following three lines
	\cleardoublepage
	\pdfbookmark{Non Disclosure}{nondisclosure}
	\thispagestyle{empty}

\chapter*{Non-disclosure Notice}

This \artderarbeit contains confidential data of \sperrvermerkUnternehmen. Any publication or reproduction of the present thesis or parts thereof it shall require the express approval of \sperrvermerkUnternehmen. Copies or transcripts of the present thesis -- including those in digital form -- shall not be allowed. Reviewers and members of the examination committee shall have access to this thesis exclusively. For determining the grade, the thesis may be made available to other persons. This non-disclosure notice shall be valid until \sperrvermerkDatum{} (for a period of 5 years). In case parts of the thesis will be or have already been published, this non-disclosure notice shall become void.

\vspace{0.5cm}

Karlsruhe, \abgabe
	
	\cleardoublepage
	\pdfbookmark{Short Description}{}
	\thispagestyle{empty}
\newgeometry{left=20mm, right=15mm, top=25mm, bottom=30mm}
\hspace{0.5cm}
\includegraphics[width=0.95\textwidth]{03_images/title_page/kitlogo_irs.pdf}

\vspace{1cm}

\centerline{\LARGE \textsf{\textbf{\artderarbeit{}\ \nummerderarbeit}}} % maybe huge,LARGE, Large

\vspace{1cm}

\begin{center}
	\Huge \textsf{\textbf{\titelderarbeit}}
\end{center}

\vspace{1cm}
\sffamily
Aufgabenstellung von Betreuer...

\vspace{9cm} % <-- Reduce space as needed


\begin{tabular}{lcl}
	Student &:& \student\addWithPreCommaSpace{\studentgrad} \\
	Supervisor &:& \betreuer\addWithPreCommaSpace{\betreuergrad} \\
	Examiner &:& Prof.\ Dr.-Ing.\ Sören Hohmann \\
	Date of hand out &:& \ausgabe \\
	Date of hand in &:& \abgabe
\end{tabular}

\vspace{0.25cm}

I declare that I wrote my \artderarbeit{} by myself and that I have followed the regulations relating to good scientific practice of the Karlsruhe Institute of Technology (KIT). I did not use any unacknowledged sources or means and I marked all references I used literally or by content.

\vspace{0.5cm}

Karlsruhe, \abgabe

\vspace{1.4cm}

\student\addWithPreCommaSpace{\studentgrad}
\rmfamily % description
	\restoregeometry % Since the description uses a different geometry
	
	\cleardoublepage
	\pdfbookmark{Acknowledgements}{acknowledgement}
	\addchap*{Acknowledgments} 
Firstly, I would like to thank ...


	
	\cleardoublepage
	\pdfbookmark{Abstract}{abstract}
	\addchap*{Abstract}

Here comes an abstract...

	
	\cleardoublepage
	\pdfbookmark{Contents}{toc}
	\tableofcontents
	
	\cleardoublepage
	\listoffigures
	
	\cleardoublepage
	\listoftables
	
	\cleardoublepage
	
\chapter*{Abbreviations and Symbols}
\label{symbol}
%\markboth{Abbreviations and Symbols}{Abbreviations and Symbols}
\addcontentsline{toc}{chapter}{Abbreviations and Symbols}

% Acronym Definitionen
\newacronym{lqr}{LQR}{Linear Quadratic Regulator}

\newacronym{mpc}{MPC}{Model Predictive Control}
\newacronym{dmpc}{DMPC}{distributed Model Predictive Control}
\newacronym[description={linear time-invariant}]{lti}{LTI}{linear time-invariant}

\printacronyms[title={Abbreviations}]
\markboth{Abbreviations and Symbols}{Abbreviations and Symbols} % Damit Kopfzeile wieder passt
%\newpage
\vspace{1.5cm}

\begin{supertabular*}{\textwidth}{ll}
% \multicolumn{2}{l}{ \Large \textbf{Abbreviations}} \\
% \\
% \hline 
% {\textbf{Abbreviation}}&{\textbf{Meaning}} \\
% \hline
% \\
% MPC & Model Predictive Control \\
% DMPC & Distributed Model Predictive Control \\
% LQR & Linear Quadratic Regulator \\
% \\
% \\
% \\
\multicolumn{2}{l}{ \Large \textbf{Latin letters}} \\
\\
\hline 
{\textbf{Symbol}}&{\textbf{Meaning}} \\
\hline
\\
$1_n$ & column vector of $n$ ones \\
$0_n$ & column vector of $n$ zeros \\
$A$ & discrete-time system matrix   \\
$A_i$ & system matrix of subsystem $i$ \\
$B$ & discrete-time input matrix\\
$B_i$ & input matrix of subsystem $i$\\
$C$ & discrete-time output matrix\\
$D$ & degree matrix of a graph \\
\\
\\
\\
\multicolumn{2}{l}{\Large \textbf{Greek letters}} \\
\\
\hline 
{\textbf{Symbol}} &	{\textbf{Meaning}} \\
\hline
\\
$\alpha$ & step size of a distributed optimization algorithm  \\
$\gamma$ & tuning factor \\
$\sigma_{\max} $ & maximum singular value \\
$\sigma_{\min} $ & minimal singular value \\
\\
\multicolumn{2}{l}{\Large \textbf{Calligraphic and other symbols}} \\
\\
\hline 
{\textbf{Symbol}} &	{\textbf{Meaning}} \\
\hline
\\
$\mathcal{E}$ & edge set of a graph \\
$\mathcal{G}$ & directed graph \\
\\
\\
\\
\multicolumn{2}{l}{ \Large \textbf{Indices, exponents and operator names}} \\
\\
\hline 
{\textbf{Symbol}} &	{\textbf{Meaning}} \\
\hline
\\
$\text{arg}$ & argument of a solution of an optimization \\
$\text{Im} A$ & image of a matrix $A$\\
$\lim$ & limit of a sequence\\
\end{supertabular*}


	\cleardoublepage
	
	\pagenumbering{arabic}
		
	%% ---------------------------------------------------------------------------
	%% Chapters
	%% ---------------------------------------------------------------------------
	
	% Add own content here. Use \cleardoublepage between chapters
	\chapter{Introduction} \label{ch:Introduction}

Here comes the introduction. Explain why the topic of your thesis is important and what is the goal of the thesis. If you need, use references like this \cite{higham2020handbook}. If you want to use acronyms, define a \textit{newacronym} in the file "sybol\_directory.tex" and use it like this: \gls{lti}. The next time you use the acronym \gls{lti} is appears authomatically in short form.  

You can also include images like this.

\begin{figure}[ht!]
    \centering
    \includegraphics[width=\textwidth]{03_images/meme.png}
    \caption{Graphical representation of the research process.}
    \label{fig:Introduction:meme}
\end{figure}


A meme can be seen in Figure~\ref*{fig:Introduction:meme}. This demonstrates that 
\begin{equation} \label{eq:Introduction:realisation}
	1 \le 2 \ \text{or} \ 1 \ge 2.
\end{equation}
In addition to \eqref{eq:Introduction:realisation}, it could be postulated that $1 \ne 2$. Further discussions are considered in the following sections and subsections.

\section{Background and Motivation} \label{sec:Introduction:Background}

\subsection{Contribution of my thesis}

text text text text text text text text text text text text text text text text text text text text text text text text text text text text text text text text text text text text text text text 
	\cleardoublepage
	\chapter{Fundamentals} \label{ch:Fundamentals}

In this chapter...
	\cleardoublepage 
	\chapter{Main Part} \label{ch:MainPart}
	
	%% ---------------------------------------------------------------------------
	%% Appendices
	%% ---------------------------------------------------------------------------
	
	\begin{appendix}
		\cleardoublepage
		\chapter{Trivial Proofs}

The following sections demonstrate the veracity of Figure~\ref{fig:Introduction:meme} and \eqref{eq:Introduction:realisation}.

\section{Central Argument}

Patrick is a genius.

\section{Motivation}

Careful observation.

\section{Listings}

Quellcodes lassen sich gut mit dem listings-Paket einbinden, so ist das Einbinden von entsprechend formatiertem Matlab-/C++/etc.-Code ohne Copy/Paste direkt aus der Quellcode-Datei möglich. Und hier noch ein Beispielhafter Matlab Quellcode:

\begin{lstlisting}[caption={Ein wahnsinnig komplizierter Quellcode}]
	%% State Space System
	asyn.SS = ss(asyn.A,asyn.B,asyn.C,asyn.D);
	
	% Infos über das System
	disp('- Informationen über das System -');
	% Ordnung des Systems
	asyn.n = rank(asyn.A);
	disp(['Ordnung des Systems n = ', num2str(asyn.n)]);
	% Polstellen
	asyn.PS = pole(asyn.SS);
	% Nullstellen
	asyn.NS = tzero(asyn.SS);
	% Beobachtbarkeit
	asyn.Ob = obsv(asyn.SS);
	if rank(asyn.Ob) == asyn.n
	disp('System vollständig Beobachtbar');
	else
	disp('System nicht vollständig Beobachtbar');
	end
	% Steuerbarkeit
	asyn.Os = ctrb(asyn.SS);
	if rank(asyn.Os) == asyn.n
	disp('System vollständig Steuerbar');
	else
	disp('System nicht vollständig Steuerbar');
	end
	
	disp('---------------------------------');
\end{lstlisting}
	\end{appendix}
	
	%% ---------------------------------------------------------------------------
	%% ---------------------------------------------------------------------------
	
	% Listings - remove if not required
	\cleardoublepage
	\lstlistoflistings % Requires the listings package
	
	% Algorithms - remove if not required
	\cleardoublepage
	\listofalgorithms % Requires the either "algorithm2e" or "algorithms" packages
	\addcontentsline{toc}{chapter}{List of algorithms}

	\cleardoublepage
	\pagenumbering{Roman}
	\printbibliography
	
\end{document}